\documentclass[UTF8, 11pt]{ctexart}

\usepackage[margin=1in]{geometry}
\usepackage{amsmath, amssymb}
\usepackage{booktabs}
\usepackage{enumitem}
\setlist{leftmargin=*}
\setlength{\parskip}{0.35em}
\setlength{\parindent}{0pt}

\title{\textbf{阶段 4：实验评估计划}\\
影响范围、运行时间与分布式近似损失}
\author{Team Members: Name 1, Name 2, Name 3}
\date{CS240: 算法设计与分析，2026 年春季学期}

\begin{document}

\maketitle

\section{运行时间分析}

比较随机基线、度数基线、朴素贪心和 CELF 的运行时间。重点观察 CELF 是否通过减少重复边际收益估计带来加速。每组实验重复运行多次并报告平均值。

\section{可扩展性评估}

在不同规模的 Erdos--Renyi 图、Barabasi--Albert 图和 Watts--Strogatz 图上评估算法表现。规模参数包括节点数 $n$、边数 $m$、预算 $k$ 和 Monte Carlo 次数 $R$。预期朴素贪心随 $n$ 和 $k$ 增长较快，而 CELF 在中等规模图上更可扩展。

\section{对比实验}

\begin{center}
\begin{tabular}{lll}
\toprule
方法 & 预期优点 & 预期缺点 \\
\midrule
随机选择 & 极快 & 影响范围低且不稳定 \\
度数中心性 & 简单、可解释 & 忽略传播重叠 \\
朴素贪心 & 近似质量高 & 运行时间高 \\
CELF & 接近贪心质量且更快 & 实现更复杂 \\
分布式变体 & 符合局部信息场景 & 可能有近似损失 \\
\bottomrule
\end{tabular}
\end{center}

主要指标包括平均影响范围、运行时间、影响力估计次数、边际收益曲线和分布式近似损失。

\section{与研究方向的联系}

实验不仅比较中心化算法，还分析局部信息和通信限制对全局目标的影响。这与分布式优化中的基本问题一致：各 agent 持有局部信息，通过有限通信逼近全局最优决策。影响力最大化提供了一个离散、子模、可实验测量的版本。

\end{document}
